\chapter{Introduction Générale}
\label{chap:intro}

\section{Contexte et Motivation}
Dans les domaines des sciences de l'ingénieur, de l'intelligence artificielle et de la recherche opérationnelle, une part prépondérante des problèmes d'optimisation réels présente des caractéristiques hautement non-linéaires, non-convexes, non-différentiables ou de grande dimension combinatoire. Face à ces fonctions objectifs complexes, les méthodes exactes classiques (telles que le simplexe, la programmation par séparation et évaluation, ou le gradient conjugué) se heurtent rapidement à l'explosion combinatoire ou à l'absence de dérivabilité analytique.

Pour résoudre ces problèmes réputés difficiles (souvent NP-difficiles), la recherche s'est orientée depuis trois décennies vers les métaheuristiques stochastiques d'inspiration biologique. Parmi celles-ci, les algorithmes d'\textbf{Intelligence en Essaim} (\textit{Swarm Intelligence}), tels que l'Optimisation par Essaim Particulaire (\textit{Particle Swarm Optimization} - PSO) introduite par Kennedy et Eberhart en 1995 \cite{kennedy1995particle}, et l'Optimisation par Loups Gris (\textit{Grey Wolf Optimizer} - GWO) conçue par Mirjalili et al. en 2014 \cite{mirjalili2014grey}, ont démontré une efficacité remarquable. Ces approches exploitent l'émergence d'une intelligence collective issue de règles de comportement simples appliquées par un ensemble d'agents décentralisés.

\section{Problématique : Le Compromis Exploration vs Exploitation}
Le cœur du succès de toute métaheuristique réside dans sa capacité à maintenir un équilibre subtil et dynamique entre deux forces antagonistes :
\begin{enumerate}[label=\textbf{\arabic*.}]
    \item \textbf{L'exploration (ou diversification) :} la capacité de l'algorithme à scruter de vastes régions inexplorées de l'espace de recherche afin de localiser les bassins d'attraction les plus prometteurs ;
    \item \textbf{L'exploitation (ou intensification) :} la capacité à concentrer l'effort de recherche localement autour des meilleures solutions déjà découvertes pour raffiner leur précision.
\end{enumerate}

Cependant, les algorithmes canoniques souffrent d'une limitation structurelle fondamentale : leurs mécanismes de recherche sont dictés par des équations déterministes et statiques. 
\begin{itemize}[label=\textbullet]
    \item Dans le \textbf{PSO canonique}, chaque particule est systématiquement attirée vers le meilleur point découvert par l'ensemble du groupe ($\mathbf{gbest}$), provoquant un phénomène d'agrégation rapide et une perte drastique de la diversité génétique de la population.
    \item Dans le \textbf{GWO canonique}, la transition exploration-exploitation repose exclusivement sur la décroissance linéaire d'un paramètre $a$ de $2$ à $0$. Ce schéma aveugle ignore totalement l'état de stagnation réel de la meute. Lorsque les leaders ($\alpha, \beta, \delta$) sont capturés par un optimum local d'un paysage accidenté (comme la fonction de Rastrigin), l'ensemble de la population converge irrémédiablement vers cette solution sous-optimale sans possibilité d'échappement.
\end{itemize}

\section{L'Apport de la Théorie des Jeux Évolutive}
Pour surmonter ces écueils sans recourir à des réglages manuels ou heuristiques ad hoc, la \textbf{Théorie des Jeux Évolutive} (\textit{Evolutionary Game Theory} - EGT) \cite{weibull1995evolutionary} fournit un cadre mathématique rigoureux et élégant. Contrairement à la théorie des jeux classique nécessitant une rationalité parfaite et le calcul coûteux d'un équilibre de Nash à chaque pas temporel, l'EGT étudie des populations d'agents dotés d'une rationalité limitée en interaction répétée.

Dans cette optique, l'essaim est modélisé comme une population de joueurs pouvant choisir parmi un portfolio de stratégies de déplacement complémentaires. Selon les performances empiriques observées (\textit{payoffs}), l'\textbf{équation du réplicateur de Taylor \& Jonker (1978)} \cite{taylor1978evolutionary} ajuste dynamiquement la représentativité de chaque stratégie : les comportements apportant une amélioration supérieure à la moyenne du groupe voient leur part croître, tandis que les stratégies inopérantes sont marginalisées.

\section{Objectifs du Projet et Structure du Rapport}
L'objectif central de ce projet d'études est de concevoir, implémenter, évaluer et analyser l'impact de l'intégration de la Théorie des Jeux Évolutive au sein de deux métaheuristiques majeures :
\begin{enumerate}[label=\textbf{(\roman*)}]
    \item L'algorithme \textbf{SDPSO} (\textit{Strategy Dynamics Particle Swarm Optimization}) \cite{liu2022strategy}, qui gère dynamiquement quatre opérateurs de vitesse (LIPS, UPSO, LDWPSO, CLPSO) sur des benchmarks continus ;
    \item L'algorithme \textbf{GT-GWO} (\textit{Game-Theoretic Grey Wolf Optimizer}), développé pour piloter quatre stratégies de chasse exclusives (\texttt{TRIO\_COOP}, \texttt{ALPHA\_EXPLOIT}, \texttt{LEVY\_EXPLORE}, \texttt{SOCIAL\_PEER}), étendu en version binaire (\textbf{BGT-GWO}) sur le problème combinatoire NP-complet du Sac à Dos 0/1.
\end{enumerate}

Le présent document est structuré en quatre chapitres :
\begin{itemize}[label=\textbullet]
    \item Le \textbf{Chapitre 1} présente l'état de l'art, les formulations mathématiques de base de PSO et GWO, les fondements de la dynamique du réplicateur en EGT, et une revue critique de la littérature.
    \item Le \textbf{Chapitre 2} détaille la conception et l'implémentation algorithmique de SDPSO, de GT-GWO continu et de son extension discrète BGT-GWO avec son opérateur de réparation glouton.
    \item Le \textbf{Chapitre 3} expose le protocole expérimental rigoureux, analyse l'ensemble des résultats chiffrés (30 à 120 répétitions), et examine les visualisations graphiques de convergence et de dynamique des stratégies.
    \item Le \textbf{Chapitre 4} propose une discussion critique sur les compromis d'exploration/exploitation, la complexité calculatoire et la portée scientifique des résultats, suivie de la conclusion générale et des perspectives.
\end{itemize}
